\documentclass[11pt]{article}

\usepackage{amsmath,amssymb}
\usepackage{xurl}
\usepackage[hidelinks]{hyperref}
\setlength{\emergencystretch}{2em}

\input{command}

% Local notation for this note.
\DeclareMathOperator{\poly}{poly}
\newcommand{\Id}{\mathrm{Id}}

\title{The Error Term of the Combined-Lines Consistency Lemma}
\author{}
\date{2026-08-30}

\begin{document}
\maketitle

\section{At a glance}

\begin{itemize}
  \item \textbf{Difficulty: easy.}  No step of the retained proof route is
  in doubt.  The obstruction is that the error function claimed in the
  statement of the combined-lines lemma of the Pauli basis analysis,
  $\poly(m^2\eps, md/q)$, is not delivered by either of the two
  derivations printed in its source proof: one route establishes
  $m \cdot \poly(\eps, md/q)$, the other would give
  $m^2 \cdot \poly(\eps, md/q)$ and rests in addition on a distributional
  decomposition that the stated sub-line lemma does not supply.  The
  substance is deciding which error form the statement should carry and
  confirming that the sole consumer tolerates it.
  \item \textbf{Key Mathlib inputs:} the
  Cauchy--Schwarz inequality for state-dependent operator distances
  (project-local, blueprint Chapter 12), mixture decompositions of
  question distributions, and elementary averaging over restricted
  line distributions; no new Mathlib surface beyond finite
  expectations and inner-product estimates already used by the
  low individual degree development.
  \item \textbf{Estimated weight:} a \textbf{statement-level error-form
  decision} plus about \textbf{one lemma interface} of checking that the
  consumer absorbs an $m$ prefactor into its constants; the three chained
  claims behind the retained route are ordinary Cauchy--Schwarz estimates,
  on the order of \textbf{a few hundred lines} of formal proof in total
  when that part of the development is reached.
  \item \textbf{Mathlib/project split:} essentially \textbf{100\%
  project-local}.  The estimates use only the project's state-dependent
  distance calculus, the sub-line distribution, and elementary
  probability on finite sets.
  \item \textbf{Key inputs:} the combined point and line measurements and
  the sub-line distribution of the Pauli basis analysis in
  \cite{Ji2020MIPStar}, and the error function of the classical soundness
  theorem imported from \cite{Ji2020LowIndividualDegree}, into which the
  corrected error form must be absorbed.  Tracked in \localissue{0002}.
\end{itemize}

\section{Key theorem forms}

Throughout, $(q,m,d)$ is an admissible parameter tuple of the Pauli basis
test, $\ket{\hat\psi}$ is the expanded state of the projective strategy
under analysis, succeeding with probability at least $1-\eps$, and all
approximations are with respect to $\ket{\hat\psi}$.  For lines
$\ell_X, \ell_Z$ in $\F_q^m$, the POVM $\{T^{\ell_X,\ell_Z}_{f_X,f_Z}\}$,
with outcomes pairs of polynomials of degree at most $md$ on the two
lines, is the two-basis combined line measurement, and
$\{\hat Q^{x,z}_{a,b}\}$ is the combined point measurement; their
consistency errors are $\delta_P = \poly(\eps, md/q)$ and
$\delta_Q = \poly(\eps)$ respectively, and
$\delta_{\mathrm{Line}} = \poly(\eps)$ is the line--point consistency
error of the strategy.  The lemma under discussion asserts that a
measurement $\{\hat Q^{\ell}_f\}$, built for each line $\ell$ in
$\F_q^{2m+2}$, is $\delta_{\mathrm{combine}}$-consistent, on average over
the line-point distribution over $\F_q^{2m+2}$, with the combined point
measurement on the extended coordinates, together with the
register-swapped counterpart.%
\footnote{In the mirror of the source consulted here,
\path{references/qpbt-paper/14_analysis_of_the_pauli_basis_test.tex}: the
lemma statement, with the error form $\poly(m^2\eps, md/q)$, is at
lines~1020--1034 (label \path{lem:qld-4-13}); the central display with
the first-route bound at lines~1134--1137
(label \path{eq:qld-combined-lines-consistency}); the three claims at
lines~1140--1239 (labels \path{claim:17-1}--\path{claim:17-3}); the
second-route conclusion $O(m^2\delta_P)$ at lines~1241--1245, preceded by
a duplicated sentence whose second copy carries a malformed reference;
the product-distribution estimate invoked there at lines~1055--1061
(label \path{eq:qld-xz-lines-restricted}); and the sub-line lemma at
lines~1063--1116 (label \path{lem:qld-sublines}).}
The comparison involves four forms of the error function
$\delta_{\mathrm{combine}}(\eps,m,d,q)$.
\begin{enumerate}
  \item \textbf{Source form.}
  $\delta_{\mathrm{combine}} = \poly(m^2\eps,\, md/q)$.  This is the form
  printed in the statement.  Neither printed derivation supports it.
  \item \textbf{Established form (first route).}
  $\delta_{\mathrm{combine}}
  = O\big(m\,(\eps^{1/4} + \delta_P^{1/4} + \delta_Q^{1/4}
  + \delta_{\mathrm{Line}}^{1/2})\big)$, a function of the form
  $m \cdot \poly(\eps,\, md/q)$.
  \item \textbf{Second-route form (conditional).}
  $\delta_{\mathrm{combine}} = O(m^2 \delta_P)
  = m^2 \cdot \poly(\eps,\, md/q)$, conditional on a joint-law
  decomposition that the sub-line lemma, as stated, does not provide.
  \item \textbf{Consumer requirement.}  The unique consumer, the
  polynomial-pair lemma, feeds $\delta_{\mathrm{combine}}$ into an error
  of the form
  $a(md)^a(\delta_{\mathrm{combine}}^b + q^{-b} + 2^{-bmd})$ imported
  from \cite{Ji2020LowIndividualDegree}; any bound
  $m^{O(1)} \cdot \poly(\eps, md/q)$ is absorbed by enlarging $a$ and
  shrinking $b$, so forms~2 and~3 both suffice downstream.
\end{enumerate}

\section{The source assertion and the two routes}

The source proof in \cite{Ji2020MIPStar} constructs $\hat Q^{\ell}_f$ by
averaging, over a distribution $D$ on tuples $(\ell, \ell_X, \ell_Z)$
provided by its sub-line lemma, the pushforward of
$T^{\ell_X,\ell_Z}$ along a combining map, and reduces the lemma to a
bound on the scalar consistency defect
\begin{equation}
  \Lambda \,=\,
  \E_{(\ell,\ell_X,\ell_Z) \sim D}\
  \E_{(x,z,\alpha,\beta) \in \ell}
  \sum_{f_X,f_Z} \bra{\hat\psi}
  \big(T^{\ell_X,\ell_Z}_{f_X,f_Z}\big)_{\mathsf{A}\mathsf{A}'} \ot
  \big(\Id - \hat{Q}^{x,z}_{f_X(x),f_Z(z)}\big)_{\mathsf{B}\mathsf{A}''}
  \ket{\hat\psi}\,,
  \tag{\path{eq:qld-combined-lines-consistency}}
\end{equation}
in which the point $(x,z,\alpha,\beta)$ is uniform on $\ell$ and the sum
ranges over pairs of line polynomials.  The proof then bounds $\Lambda$
twice, by different routes and with different results, and the printed
record is internally inconsistent: the display introducing $\Lambda$
asserts the first-route bound, while the closing lines assert the
second-route bound as the operative one; the lemma statement claims a
third form matching neither.

\paragraph{The first route.}
Three claims are chained.  The combined point measurement
$\hat Q^{x,z}$ is replaced by the ordered product of the two point
measurements at cost $\delta_Q^{1/2}$; the $X$-point factor is removed at
cost $m\,\delta_{\mathrm{Line}}^{1/2}$; and the remaining $Z$-point
correlation is shown to be within
$\sqrt{m}\,(\delta_P^{1/4} + \delta_Q^{1/4} + \eps^{1/4})$ of one.  The
factors of $m$ enter through the mixture structure of the sub-line
distribution: each restricted line distribution carries mixture weight
$1/(2m)$ in the line-point distribution, so an average bounded by
$\delta$ over the full distribution is bounded only by $2m\delta$ over a
restricted component and by $4m^2\delta$ over a product of two
components, and the Cauchy--Schwarz steps halve the exponents.  Chaining
the claims gives
\begin{equation*}
  \Lambda \,=\, O\big(m\,(\eps^{1/4} + \delta_P^{1/4} + \delta_Q^{1/4}
  + \delta_{\mathrm{Line}}^{1/2})\big)
  \,=\, m \cdot \poly(\eps,\, md/q)\,.
\end{equation*}
One intermediate assertion needs a harmless correction: the final line
of the proof of the second claim asserts the bound
$O(m^2 \delta_{\mathrm{Line}})$ for a consistency term that is close to
one, rather than for its deficit.  Read as a bound on the deficit ---
which is what the surrounding Cauchy--Schwarz argument requires and what
the invoked line--point consistency provides after the mixture
inflation --- the claim and its use are correct, with the square root of
the deficit giving the stated cost $m\,\delta_{\mathrm{Line}}^{1/2}$.

\paragraph{The second route.}
After the chain, the source restarts: it asserts that, by Property~2 of
its sub-line lemma, $\Lambda$ is a mixture, over types
$\mathrm{v}_1, \mathrm{v}_2 \in \{\mathrm{ALine}, \mathrm{DLine}\}$ and
indices $i, j \in \{1, \dots, m\}$, of terms
\begin{equation*}
  \E_{(\ell_X, x) \sim D_{\mathrm{v}_1,i}}\
  \E_{(\ell_Z, z) \sim D_{\mathrm{v}_2,j}}
  \sum_{f_X,f_Z} \bra{\hat\psi}
  \big(T^{\ell_X,\ell_Z}_{f_X,f_Z}\big)_{\mathsf{A}\mathsf{A}'} \ot
  \big(\Id - \hat{Q}^{x,z}_{f_X(x),f_Z(z)}\big)_{\mathsf{B}\mathsf{A}''}
  \ket{\hat\psi}
\end{equation*}
over two \emph{independently} sampled restricted line-point pairs,
bounds each term by $O(m^2\delta_P)$ using the restricted-distribution
consequence of the two-basis line lemma, and concludes
$\Lambda = O(m^2\delta_P)$.  The invoked decomposition concerns the
\emph{joint} law of the pair $((\ell_X, x), (\ell_Z, z))$, where $x$ and
$z$ are the two blocks of the single point sampled on $\ell$.  Property~2
of the sub-line lemma does not assert this joint law: it asserts the
component-wise product law of the pair of \emph{lines}
$(\ell_X, \ell_Z)$, and the conditional uniformity of each point on its
line separately --- equivalently, the laws of the two pairs
$((\ell_X, x), (\ell_Z, z'))$ and $((\ell_X, x'), (\ell_Z, z))$ in which
the \emph{other} point is resampled independently.  These marginal
statements are exactly what the first route's claims consume, each
involving one point at a time; the joint decomposition used by the
second route is a strictly stronger property and is not established by
the lemma it is attributed to.

\section{Why neither route yields the stated form}

The stated error form is $\poly(m^2\eps, md/q)$, a function of the two
quantities $m^2\eps$ and $md/q$ that vanishes as both vanish.  The
first-route bound is not dominated by any such function: along the
scaling $\eps = m^{-4}$ with $q/(md) \to \infty$, both $m^2\eps = m^{-2}$
and $md/q$ tend to $0$, while the term $m\,\eps^{1/4} = 1$ does not.
Hence the established derivation supports no bound of the claimed form.
The second-route bound, even granting the missing decomposition, carries
the separate prefactor $m^2$: writing $\delta_P = \poly(\eps, md/q)$
with some fixed positive exponents, the contribution
$m^2 (md/q)^{c}$ is not dominated by any vanishing function of $m^2\eps$
and $md/q$, as the scaling $\eps = 0$, $md/q = m^{-2/c}$ shows, and the
$\eps$-contribution fails similarly for sublinear exponents.  Neither
observation asserts that the \emph{conclusion} of the lemma is false for
some error function of the claimed form; the defect is that the printed
proof does not deliver it.

For orientation, the algebraic identity
$m\,\eps^{1/2} = (m^2\eps)^{1/2}$ shows what a derivation of the claimed
$\eps$-dependence would need: a chain that loses only square roots in the
$\eps$-type quantities \emph{after} the restricted-distribution
inflation would produce terms of the shape $(m^2\eps)^{1/2}$.  The
printed chain loses fourth roots, through two nested Cauchy--Schwarz
steps applied to already square-rooted quantities, and its
$md/q$-dependence enters as $m\,\delta_P^{1/4}$, which the claimed form
cannot absorb.

\section{Candidate repairs}

The following are candidate directions, with their open obligations; none
of the routes below is claimed as established.

\paragraph{Correcting the error form (recommended direction).}
Adopt $\delta_{\mathrm{combine}} = m \cdot \poly(\eps, md/q)$ ---
concretely, the first-route bound --- as the operative error form.  Open
obligations: (i)~verify that every consumer tolerates the prefactor; in
the source and in the blueprint the only consumer is the polynomial-pair
lemma, whose imported error
$a(md)^a(\delta_{\mathrm{combine}}^b + q^{-b} + 2^{-bmd})$ absorbs
$m^{O(1)}$ prefactors by enlarging $a$ and shrinking $b$, and this
absorption must be carried out with explicit constants when the lemma is
formalized; (ii)~record the correction against the source-labelled
statement in accordance with the proof-gap protocol, rather than
silently substituting the corrected form --- the blueprint currently
retains the source statement and carries the corrected bound in the
proof; (iii)~confirm the three chained claims at the required level of
detail, including the corrected deficit reading of the second claim.

\paragraph{Repairing the second route.}
Strengthen Property~2 of the sub-line lemma to the joint product law: in
each mixture component, the pairs $(\ell_X, x)$ and $(\ell_Z, z)$ are
independent with the stated restricted-distribution laws.  The sampling
procedure makes this plausible, since in each component the two pairs
are built from disjoint collections of the independent ingredients (the
two blocks of the uniform point, the direction blocks, and the fresh
auxiliary choices).  Open obligations: (i)~state and prove the
strengthened property for every component, including lines whose
direction blocks vanish and degenerate to singletons; (ii)~re-derive the
mixture decomposition of $\Lambda$ from the strengthened property;
(iii)~note that even then the route delivers only
$O(m^2\delta_P) = m^2 \cdot \poly(\eps, md/q)$, so this direction removes
the unsupported step of the printed proof but still requires the
error-form decision of the previous direction, with $m^2$ in place of
$m$.

\paragraph{Establishing the claimed form.}
Re-derive the chain so that every factor of $m$ enters jointly with
$\eps$ as $m^2\eps$, in the sense of the orientation identity above.
Open obligations: (i)~find a chaining that avoids one of the two nested
square-root losses, or prove restricted-distribution versions of the
underlying consistency lemmas directly at their inflated errors;
(ii)~handle the $md/q$-dependence, since a residual term
$m (md/q)^{c}$ is incompatible with the claimed form for the current
exponents.  No derivation along these lines is known here, and the
direction may be infeasible; in that case the corrected error form of
the first direction stands.

\section{Comparison with the blueprint}

The blueprint \cite{MIPStarREBlueprint} keeps the combined-lines lemma in
its source shape, error form $\poly(m^2\eps, md/q)$ included, and marks
in the proof that this form remains unproved: the proof derives the
first route only, records the established bound
$m \cdot \poly(\eps, md/q)$, and uses that bound in the sequel.  The
three claims are stated as separate lemmas with the corrected deficit
assertion for the second; the sub-line lemma states Property~2 at the
level of the two marginal laws and expressly disclaims the joint
conditional law of the two points, so the blueprint statement does not
license the second route; and the polynomial-pair lemma's proof absorbs
the $m$ prefactor into the constants of the imported soundness error.  A
chapter remark summarizes the discrepancy and cites this note.%
\footnote{The blueprint statements are \path{lem:qld-4-13},
\path{lem:qld-sublines}, \path{lem:claim-17-1}--\path{lem:claim-17-3},
\path{lem:qld-4-7}, and the remark
\path{rem:qld-4-13-source-defects}, all in
\path{blueprint/src/chapter/ch15_qpbt_combining.tex}.  No Lean
declaration is currently linked to any of them.}

\section{Conclusion}

There is a proof-level gap in the cited lemma, together with an
intermediate inconsistency: the printed proof asserts two different
bounds for the same consistency defect, the first-route bound
$O(m(\eps^{1/4} + \delta_P^{1/4} + \delta_Q^{1/4} +
\delta_{\mathrm{Line}}^{1/2}))$ and the second-route bound
$O(m^2\delta_P)$, the second resting on a joint-law decomposition that
the sub-line lemma as stated does not supply; and the error form claimed
in the lemma statement, $\poly(m^2\eps, md/q)$, follows from neither.
The first route is sound after a harmless deficit correction in its
second claim and establishes $m \cdot \poly(\eps, md/q)$.  The
discrepancy is harmless downstream, since the sole consumer absorbs
polynomial-in-$m$ prefactors into its constants.  The recommended
resolution is to treat the corrected error form as operative while the
source-labelled statement retains its printed form with the missing step
explicitly tracked; strengthening the sub-line lemma to license the
second route is optional, and no route to the claimed form is currently
known.  We do not assert that the claimed form is false.

\bibliographystyle{alpha}
\bibliography{references}

\end{document}
